% !TEX program = xelatex
% AI-effects 프로젝트 --- AI 점유율 지표 A~C(paper 27)와 노출도 지표 3개의 상관과 상위 순위 비교.
% 표 본문: results/table/28_corr.tex, 28_overlap.tex, 28_top_fs.tex, 28_top_expo.tex
%          (code/04_analysis_kor_share_vs_exposure.do가 생성)
% 전체 자료: results/table/28_kor_share_vs_exposure.xlsx
% 관련: paper 25(활용도와 노출도), paper 27(점유율 지표)
% 컴파일: xelatex (kotex 필요)
\documentclass[11pt]{article}

\usepackage{fontspec}
\usepackage{kotex}
\usepackage[a4paper,margin=2.0cm]{geometry}
\usepackage{longtable}
\usepackage{booktabs}
\usepackage{array}
\usepackage{amsmath}
\usepackage{xcolor}
\usepackage{hyperref}
\usepackage{enumitem}

\hypersetup{colorlinks=true,linkcolor=blue!50!black,urlcolor=blue!50!black}
\newcolumntype{L}[1]{>{\raggedright\arraybackslash}p{#1}}
\newcolumntype{R}[1]{>{\raggedleft\arraybackslash}p{#1}}
\setlength{\LTleft}{0pt}
\setlength{\LTright}{0pt}
\renewcommand{\arraystretch}{1.15}
\setlength{\tabcolsep}{3.5pt}

\title{AI 점유율 지표와 노출도 지표의 비교\\[2mm]
\large --- 지표 A$\sim$C와 Felten AIOE, Eloundou GPT-4 $\beta$, Observed exposure ---}
\author{AI-effects 프로젝트}
\date{\today}

\begin{document}
\maketitle

\begin{abstract}
\noindent paper 27의 AI 점유율 지표 A$\sim$C를 기존 노출도 지표 3개와 KSCO 8차 소분류 129개에서 비교했다. 노출도는 직업의 강도 지표이고 점유율은 강도에 규모가 곱해진 지표라서, 규모를 맞춘 ``노출도 기반 점유율''(취업자 $\times$ 노출도의 점유율)도 함께 보았다. 주요 결과는 다음과 같다.
\begin{itemize}[leftmargin=*]
  \item 점유율 지표와 노출도 지표의 순위상관은 0.60$\sim$0.70이다. Observed exposure와 Eloundou $\beta$가 0.68$\sim$0.70으로 가장 높고, Felten이 0.60$\sim$0.63으로 가장 낮다.
  \item 규모를 맞춘 노출도 기반 점유율과 지표 C의 순위상관은 0.50(Observed exposure), 0.22(Eloundou $\beta$)로 오히려 낮다. 노출도로 보면 취업자가 많은 사무$\cdot$판매직(회계 및 경리 사무원, 기획 및 마케팅 사무원, 매장 판매 종사자)이 상위지만, 실제 사용 점유율은 이들 직업에서 0.1\% 안팎이다.
  \item 상위 30개 겹침도 같은 그림이다. 지표 C 상위 30개 중 노출도 상위 30개와 겹치는 것은 15$\sim$17개, 노출도 기반 점유율 상위 30개와는 8$\sim$11개다.
\end{itemize}
\end{abstract}

\section{비교 방법}

\begin{itemize}[leftmargin=*]
  \item \textbf{대상}: 지표 A$\sim$C와 노출도 3개가 모두 있는 소분류 129개(paper 25와 같은 표본).
  \item \textbf{점유율 지표}(paper 27): A 대화 점유율, B 절감 시간 점유율, C 취업자 환산 절감 점유율. $K$와 무관하고 직업 규모로 나누지 않는다.
  \item \textbf{노출도 지표}: Observed exposure(Anthropic), Felten 언어모델 AIOE, Eloundou GPT-4 $\beta$. SOC 값을 연계표 $w^{\text{mean}}$으로 평균한 소분류 값이다(paper 23). 모두 직업 규모와 무관한 강도 지표다.
  \item \textbf{노출도 기반 점유율}(작업자 판단): $\mathrm{XS}^x_j = E_j x_j / \sum_k E_k x_k$. 여기서 $E_j$는 취업자 수(2022년 상반기), $x$는 Observed exposure 또는 Eloundou $\beta$다. ``노출도만으로 본다면 영향받는 일자리가 어느 직업에 몰리는가''를 나타내며, 지표 C와 같은 단위(점유율)다. Felten AIOE는 표준화 점수(음수 포함)라 만들지 않았다.
  \item \textbf{통계}: Spearman 순위상관(주), Pearson 상관(엑셀), 상위 30개 겹침 수, 지표 C 상위 20개의 각 지표 순위, 노출도별 상위 10개의 지표 C 순위.
\end{itemize}

\section{상관관계}

\begin{longtable}{L{4.6cm}R{1.4cm}R{1.4cm}R{1.4cm}R{1.4cm}R{1.4cm}R{1.4cm}R{1.4cm}R{1.4cm}}
\caption{Spearman 순위상관(소분류 129개)}\label{tab:corr}\\
\toprule
 & (1) & (2) & (3) & (4) & (5) & (6) & (7) & (8) \\
\midrule
\endfirsthead
% 04_analysis_kor_share_vs_exposure.do가 생성. 직접 고치지 말 것.
(1) A 대화 점유율 & -- &  &  &  &  &  &  &  \\
(2) B 절감 시간 점유율 & 0.96 & -- &  &  &  &  &  &  \\
(3) C 취업자 환산 점유율 & 0.96 & 1.00 & -- &  &  &  &  &  \\
(4) Observed exposure & 0.70 & 0.69 & 0.69 & -- &  &  &  &  \\
(5) Felten LM AIOE & 0.60 & 0.63 & 0.63 & 0.81 & -- &  &  &  \\
(6) Eloundou GPT-4 $\beta$ & 0.68 & 0.70 & 0.69 & 0.84 & 0.85 & -- &  &  \\
(7) 노출도 점유율(OE) & 0.55 & 0.50 & 0.50 & 0.76 & 0.54 & 0.57 & -- &  \\
(8) 노출도 점유율($\beta$) & 0.29 & 0.21 & 0.22 & 0.30 & 0.20 & 0.25 & 0.78 & -- \\

\bottomrule
\end{longtable}

\begin{itemize}[leftmargin=*]
  \item 점유율 지표와 노출도의 순위상관은 0.60$\sim$0.70이다. paper 25의 강도 지표(활용도 1$\cdot$2)와 노출도의 순위상관(0.58$\sim$0.66)과 비슷하거나 조금 높다. 노출도 지표끼리의 상관(0.81$\sim$0.85)보다는 낮다. 실제 사용은 노출도와 같은 방향이지만 다른 정보를 담는다.
  \item 세 노출도 가운데 Felten AIOE가 점유율과 가장 덜 맞는다(0.60$\sim$0.63). Felten 상위에는 통신 관련 판매, 종교 관련 종사자, 영양사처럼 실제 사용이 적은 직업이 있다(표~\ref{tab:topexpo}).
  \item \textbf{규모를 맞추면 상관이 더 낮아진다.} 노출도 기반 점유율과 지표 C의 순위상관은 Observed exposure로 0.50, Eloundou $\beta$로 0.22다. 노출도 기반 점유율 상위는 회계 및 경리 사무원(331, Eloundou 기준 7.8\%), 기획 및 마케팅 사무원(311, 6.9\%), 매장 판매 종사자(521), 자재 및 생산관리 사무원(321)이다. 이들의 지표 C 점유율은 0.06$\sim$0.24\%다. 노출도는 취업자가 많은 중간 숙련 사무$\cdot$판매직에 일자리 영향이 몰린다고 보지만, 실제 AI 사용으로 본 영향은 전문가 직업에 몰린다(paper 27의 대분류 표).
\end{itemize}

\section{상위 순위 비교}

\begin{longtable}{L{4.6cm}R{1.4cm}R{1.4cm}R{1.4cm}R{1.4cm}R{1.4cm}R{1.4cm}R{1.4cm}R{1.4cm}}
\caption{상위 30개 겹침 수}\label{tab:overlap}\\
\toprule
 & (1) & (2) & (3) & (4) & (5) & (6) & (7) & (8) \\
\midrule
\endfirsthead
% 04_analysis_kor_share_vs_exposure.do가 생성. 직접 고치지 말 것.
(1) A 대화 점유율 & -- &  &  &  &  &  &  &  \\
(2) B 절감 시간 점유율 & 26 & -- &  &  &  &  &  &  \\
(3) C 취업자 환산 점유율 & 26 & 29 & -- &  &  &  &  &  \\
(4) Observed exposure & 18 & 17 & 17 & -- &  &  &  &  \\
(5) Felten LM AIOE & 15 & 15 & 15 & 20 & -- &  &  &  \\
(6) Eloundou GPT-4 $\beta$ & 15 & 15 & 15 & 24 & 19 & -- &  &  \\
(7) 노출도 점유율(OE) & 13 & 11 & 11 & 16 & 15 & 17 & -- &  \\
(8) 노출도 점유율($\beta$) & 10 & 8 & 8 & 12 & 10 & 13 & 23 & -- \\

\bottomrule
\end{longtable}

\begin{longtable}{R{0.6cm}L{4.4cm}R{1.2cm}R{1.0cm}R{1.0cm}R{1.2cm}R{1.2cm}R{1.2cm}R{1.4cm}R{1.4cm}}
\caption{지표 C 상위 20개 소분류의 각 지표 순위(129개 중)}\label{tab:topfs}\\
\toprule
C 순위 & 소분류 & C(\%) & A & B & OE & Felten & $\beta$ & 노출 점유율(OE) & 노출 점유율($\beta$) \\
\midrule
\endfirsthead
\toprule
C 순위 & 소분류 & C(\%) & A & B & OE & Felten & $\beta$ & 노출 점유율(OE) & 노출 점유율($\beta$) \\
\midrule
\endhead
% 04_analysis_kor_share_vs_exposure.do가 생성. 직접 고치지 말 것.
1 & 222 컴퓨터 시스템 및 소프트웨어 전문가 & 11.56 & 2 & 1 & 2 & 40 & 3 & 3 & 5 \\
2 & 291 작가 및 언론 관련 전문가 & 11.08 & 1 & 2 & 10 & 12 & 9 & 21 & 30 \\
3 & 372 비서 및 사무 지원 종사자 & 5.81 & 4 & 4 & 16 & 45 & 14 & 10 & 10 \\
4 & 225 정보 시스템 및 웹 운영자 & 5.13 & 3 & 3 & 5 & 47 & 7 & 19 & 38 \\
5 & 283 상품 기획‧홍보 및 조사 전문가 & 5.06 & 5 & 5 & 6 & 17 & 12 & 23 & 47 \\
6 & 332 통계‧데이터 관련 사무원 & 4.20 & 9 & 6 & 9 & 11 & 5 & 55 & 93 \\
7 & 213 생명 및 자연과학 관련 시험원 & 3.51 & 8 & 8 & 40 & 83 & 52 & 86 & 107 \\
8 & 223 네트워크 및 정보 보안 전문가 & 3.51 & 6 & 7 & 3 & 42 & 10 & 31 & 68 \\
9 & 293 연극‧영화 및 영상 전문가 & 2.91 & 11 & 10 & 56 & 66 & 64 & 41 & 42 \\
10 & 211 생명 및 자연과학 관련 전문가 & 2.78 & 12 & 9 & 48 & 43 & 31 & 45 & 50 \\
11 & 264 문리‧기술 및 예능 강사 & 2.74 & 14 & 13 & 44 & 39 & 60 & 12 & 8 \\
12 & 292 학예사‧사서 및 기록물 관리사 & 2.37 & 7 & 12 & 25 & 60 & 36 & 53 & 86 \\
13 & 269 기타 교육 전문가 & 2.33 & 16 & 14 & 18 & 5 & 46 & 38 & 75 \\
14 & 282 금융 및 보험 전문가 & 2.30 & 13 & 11 & 11 & 15 & 22 & 32 & 64 \\
15 & 340 금융 사무 종사자 & 1.73 & 10 & 16 & 36 & 19 & 16 & 15 & 15 \\
16 & 224 데이터 전문가 & 1.72 & 21 & 15 & 1 & 28 & 1 & 35 & 70 \\
17 & 235 기계‧로봇공학 기술자 및 시험원 & 1.59 & 15 & 17 & 74 & 56 & 40 & 58 & 37 \\
18 & 122 마케팅 및 광고‧홍보 관리자 & 1.44 & 25 & 18 & 24 & 9 & 28 & 84 & 113 \\
19 & 261 대학교수 및 강사 & 1.33 & 26 & 24 & 39 & 2 & 34 & 27 & 39 \\
20 & 131 연구‧교육 및 법률 관련 관리자 & 1.28 & 17 & 19 & 57 & 29 & 43 & 70 & 69 \\

\bottomrule
\end{longtable}

\begin{longtable}{L{2.6cm}R{0.8cm}L{5.4cm}R{1.4cm}R{1.6cm}R{1.6cm}R{1.4cm}}
\caption{노출도 지표별 상위 10개 소분류와 지표 C}\label{tab:topexpo}\\
\toprule
노출도 지표 & 순위 & 소분류 & 노출도 값 & 취업자(천 명) & C 점유율(\%) & C 순위 \\
\midrule
\endfirsthead
\toprule
노출도 지표 & 순위 & 소분류 & 노출도 값 & 취업자(천 명) & C 점유율(\%) & C 순위 \\
\midrule
\endhead
% 04_11_analysis_kor_share_vs_exposure.do가 생성. 직접 고치지 말 것.
Observed exposure & 1 & 224 데이터 전문가 & 0.46 & 25 & 1.72 & 16 \\
 & 2 & 222 컴퓨터 시스템 및 소프트웨어 전문가 & 0.40 & 399 & 11.56 & 1 \\
 & 3 & 223 네트워크 및 정보 보안 전문가 & 0.39 & 33 & 3.51 & 8 \\
 & 4 & 362 고객 상담 및 모니터 요원 & 0.36 & 228 & 0.28 & 53 \\
 & 5 & 225 정보 시스템 및 웹 운영자 & 0.35 & 79 & 5.13 & 4 \\
 & 6 & 283 상품 기획‧홍보 및 조사 전문가 & 0.35 & 62 & 5.06 & 5 \\
 & 7 & 312 인사 및 교육·훈련 사무원 & 0.31 & 203 & 0.37 & 46 \\
 & 8 & 361 여행‧안내 및 접수 사무원 & 0.31 & 119 & 0.46 & 40 \\
 & 9 & 332 통계‧데이터 관련 사무원 & 0.30 & 16 & 4.20 & 6 \\
 & 10 & 291 작가 및 언론 관련 전문가 & 0.28 & 96 & 11.08 & 2 \\
\midrule
Felten LM AIOE & 1 & 531 통신 관련 판매 종사자 & 1.86 & 255 & 0.11 & 72 \\
 & 2 & 261 대학교수 및 강사 & 1.52 & 113 & 1.33 & 19 \\
 & 3 & 271 법률 전문가 & 1.42 & 33 & 1.19 & 23 \\
 & 4 & 254 종교 관련 종사자 & 1.41 & 116 & 0.10 & 73 \\
 & 5 & 269 기타 교육 전문가 & 1.40 & 46 & 2.33 & 13 \\
 & 6 & 212 인문 및 사회과학 전문가 & 1.37 & 22 & 1.15 & 24 \\
 & 7 & 245 영양사 & 1.36 & 40 & 0.02 & 98 \\
 & 8 & 312 인사 및 교육·훈련 사무원 & 1.31 & 203 & 0.37 & 46 \\
 & 9 & 122 마케팅 및 광고‧홍보 관리자 & 1.27 & 9 & 1.44 & 18 \\
 & 10 & 532 방문 및 노점 판매 종사자 & 1.27 & 85 & 0.02 & 104 \\
\midrule
Eloundou GPT-4 $\beta$ & 1 & 224 데이터 전문가 & 0.92 & 25 & 1.72 & 16 \\
 & 2 & 311 기획 및 마케팅 사무원 & 0.84 & 707 & 0.06 & 84 \\
 & 3 & 222 컴퓨터 시스템 및 소프트웨어 전문가 & 0.81 & 399 & 11.56 & 1 \\
 & 4 & 331 회계 및 경리 사무원 & 0.80 & 826 & 0.23 & 57 \\
 & 5 & 332 통계‧데이터 관련 사무원 & 0.79 & 16 & 4.20 & 6 \\
 & 6 & 362 고객 상담 및 모니터 요원 & 0.75 & 228 & 0.28 & 53 \\
 & 7 & 225 정보 시스템 및 웹 운영자 & 0.72 & 79 & 5.13 & 4 \\
 & 8 & 531 통신 관련 판매 종사자 & 0.71 & 255 & 0.11 & 72 \\
 & 9 & 291 작가 및 언론 관련 전문가 & 0.71 & 96 & 11.08 & 2 \\
 & 10 & 223 네트워크 및 정보 보안 전문가 & 0.70 & 33 & 3.51 & 8 \\

\bottomrule
\end{longtable}

\begin{itemize}[leftmargin=*]
  \item \textbf{두 지표가 모두 높은 직업}: 222 컴퓨터 시스템 및 소프트웨어 전문가, 225 정보 시스템 및 웹 운영자, 223 네트워크 및 정보 보안 전문가, 332 통계$\cdot$데이터 관련 사무원, 291 작가 및 언론 관련 전문가는 노출도(Observed exposure, Eloundou $\beta$) 상위 10과 지표 C 상위 10에 함께 든다. 노출된 과업이 실제로도 많이 쓰이는 정보통신$\cdot$글쓰기 직업이다.
  \item \textbf{사용이 노출도보다 큰 직업}: 213 생명 및 자연과학 관련 시험원(C 7위, 노출도 40$\sim$83위), 293 연극$\cdot$영화 및 영상 전문가(C 9위, 56$\sim$66위), 211 생명 및 자연과학 관련 전문가(C 10위, 31$\sim$48위). 연계표 배분의 영향(Data Scientists 등이 들어옴, paper 26)과 노출도 지표가 포착하지 못한 사용이 섞여 있다.
  \item \textbf{노출도가 사용보다 큰 직업}: 362 고객 상담 및 모니터 요원(Observed exposure 4위, C 53위), 311 기획 및 마케팅 사무원($\beta$ 2위, C 84위), 331 회계 및 경리 사무원($\beta$ 4위, C 57위), 531 통신 관련 판매 종사자(Felten 1위, C 72위). 취업자가 많은 사무$\cdot$판매직이 많다.
  \item 노출도 상위 10개의 지표 C 순위 중앙값은 Observed exposure 7위, Eloundou $\beta$ 12위, Felten 35위다. 지표 C 상위 30에 드는 것은 각각 7개, 6개, 5개다.
\end{itemize}

\section{해석상 주의}

\begin{enumerate}[leftmargin=*]
  \item \textbf{지표의 성격 차이}: 노출도는 ``AI가 할 수 있는 과업의 비중''(강도), 점유율은 ``한국 전체 사용 중 그 직업의 몫''(규모)이다. 노출도 기반 점유율은 둘을 맞추려는 보조 지표일 뿐 표준 지표가 아니다.
  \item \textbf{사무직의 낮은 사용 점유율}: AEI는 대화를 과업이 속한 SOC 직업으로 분류한다. 회계$\cdot$마케팅 관련 대화가 Accountants, Market Research Analysts 같은 전문가 SOC로 분류되면, 연계표를 거쳐 사무원이 아닌 전문가 소분류(예: 283 상품 기획$\cdot$홍보 및 조사 전문가)로 간다. 사무직의 낮은 점유율 일부는 실제 사용이 적어서가 아니라 이 분류 경로 때문일 수 있다.
  \item \textbf{같은 연계 경로}: 노출도와 점유율 모두 SOC 값을 같은 연계표로 옮겼다. 다만 노출도는 $w^{\text{mean}}$ 평균, 점유율은 $w^{\text{alloc}}$ 배분이라 연계 경로의 영향이 다르게 나타난다.
  \item \textbf{Observed exposure}: Massenkoff and McCrory(2026)가 Eloundou의 이론적 노출등급과 AEI 실제 사용 자료를 결합해 만든 지표다. 점유율과 출처(AEI)가 일부 겹치므로 상관이 높게 나오는 것을 독립적인 확인으로 보기 어렵다. Eloundou $\beta$와의 상관(0.84)이 높은 것도 같은 이유다.
\end{enumerate}

\section{산출 파일}

\begin{itemize}[leftmargin=*]
  \item \texttt{code/04\_analysis\_kor\_share\_vs\_exposure.do}: 비교 지표, 상관, 순위 표를 만든다.
  \item \texttt{results/table/28\_kor\_share\_vs\_exposure.xlsx}: 시트 \texttt{ksco3}(129개 소분류의 8개 지표와 순위), \texttt{spearman}, \texttt{pearson}, \texttt{overlap30}(8$\times$8 행렬), \texttt{top\_fs}, \texttt{top\_expo}.
  \item \texttt{results/table/28\_corr.tex}, \texttt{28\_overlap.tex}, \texttt{28\_top\_fs.tex}, \texttt{28\_top\_expo.tex}: 표~\ref{tab:corr}--\ref{tab:topexpo}.
\end{itemize}

\end{document}
